% =====================================================================
% notation.tex -- table of notation (included in Section 2).
% Draft of step 6.2 (completed for Sections 2-3 in step 6.3), from model.tex and the notes (sections 4.1,
% 4.14-4.26). Complete in 6.3-6.6 as symbols enter the text; prune in 6.8.
%
% Decisions of step 6.2 (clashes resolved):
%   - Periods are t (s) and tau (units of L/c_r), like t_a and tau_a:
%     t_1, tau_1 for the loop, t_s, tau_s for a fixed-free strand,
%     t_{B1}, t_{A1} for modes followed by strand. T stays for tensions
%     (T_0, T_d, T_t, and T_1, T_2 at the drive as in DIN 22101).
%   - Transit time of strand B: tau_B = 1 - xi + gamma (was t_B).
%   - Drive factor written e^{mu theta}; E is only Young's modulus.
%     Wrap angle theta (DIN 22101 uses alpha, taken here by the material
%     coupling factor); friction coefficient mu without subscript, while
%     line densities always carry one (mu_r, mu_c).
%   - Inclination of the belt path delta(s), as in DIN 22101 (model.tex
%     used theta(s)); carriage inclination delta_c (model.tex used kappa
%     for its sine, which clashed with the wavenumbers kappa_j).
%   - Weighing line density m' (DIN 22101); m is kept for masses
%     (modal mass m_k, strand masses m_A, m_B, drive mass ratio m_d).
%   - Rigging ratio lambda (model.tex used i, the imaginary unit).
%   - T_t/F_U gets no symbol (the notes of step 5.4 used tau).
% =====================================================================
\begin{table}[p]
\centering
\caption{Notation. Dimensionless quantities are scaled with the strand length $L$,
the return-strand transit time $L/c_r$, the peak drive acceleration $a_m$ and the
inertial force $\mu_rLa_m$.}
\label{tab:notation}
\footnotesize
\begin{tabular}{@{}l>{\raggedright\arraybackslash}p{0.72\linewidth}@{}}
\toprule
Symbol & Meaning\\
\midrule
\multicolumn{2}{@{}l}{\emph{Geometry}}\\
$L$ & path length of each strand (carry and return)\\
$\sigma$ & geometric coordinate from the head pulley in the direction of travel; return strand $0\le\sigma<L$\\
$\sigma_d$, $\sigma_t$ & positions of the drive and take-up pulleys\\
$s$, $x=s/L$ & loop coordinate from the drive exit; $s=0$ and $s=2L$ are the faces of the drive pulley\\
$s_t$, $\xi=s_t/L$ & take-up position from the drive exit\\
$\ell_1$, $\ell_2$ & return belt between head and drive, and between take-up and tail (units of $L$)\\
$\delta$ & local inclination of the belt path ($\delta>0$ uphill)\\
\midrule
\multicolumn{2}{@{}l}{\emph{Belt, drive and take-up}}\\
$EA$ & axial stiffness of the belt ($E$, Young's modulus)\\
$\mu_r$, $\mu_c$ & inertial line densities (belt, reduced idler mass, coupled material)\\
$\mw_r$, $\mw_c$ & line densities that contribute to weight (belt and material)\\
$c_r$, $c_c$ & wave speeds, $c_j=\sqrt{EA/\mu_j}$; $Z_r=\mu_rc_r$, impedance of the return strand\\
$\alpha$ & material coupling factor \citep{lodewijks2002}\\
$t_v$ & Kelvin--Voigt retardation time\\
$\Mtu$; $M_w$, $M_c$ & take-up mass referred to the carriage; counterweight and carriage masses\\
$\nstr$ & belt strands that carry the take-up pulley\\
$\lambda$ & rigging ratio (counterweight travel per unit carriage travel; $\lambda=1$ when hung directly)\\
$\delta_c$ & inclination of the take-up carriage\\
$\Tt$ & static belt tension at the take-up\\
$T_1$, $T_2$ & belt tensions at the drive entry and exit (tight and slack side)\\
$\Wr=g\mu_rL$ & weight of the return strand\\
$F_U$ & running peripheral force at the drive (DIN 22101)\\
$\egrip$ & drive factor (friction coefficient $\mu$ between belt and pulley, wrap angle $\theta$)\\
$m_d$, $\hat c$ & reduced drive mass over $\mu_rL$ and slip damping over $\mu_rc_r$ (drive without speed control)\\
\midrule
\multicolumn{2}{@{}l}{\emph{Response}}\\
$w$, $U$ & elastic displacement of the belt; rigid motion imposed by the drive\\
$y$, $y_c$ & take-up displacement (two-strand equivalent, positive when the loop lengthens); carriage displacement\\
$T_0$, $T_d$ & static and dynamic belt tension; $\hat T=T_d/(\mu_rLa_m)$\\
$V$, $a$, $a_m$ & drive velocity, acceleration and peak acceleration\\
$t_a$, $\tau_a$ & start (acceleration) time; $\tau_a=c_rt_a/L$\\
$r$, $\varphi$ & motion resistance per unit length; its onset function\\
\midrule
\multicolumn{2}{@{}l}{\emph{Dimensionless parameters and modal quantities}}\\
$\gamma=c_r/c_c$ & wave-speed ratio, $\gamma^2=\mu_c/\mu_r$\\
$\beta$ & take-up mass ratio, $4\Mtu/(\nstr^2\mu_rL)=(4/\nstr)\lambda\,\Tt/\Wr$ (carriage mass neglected)\\
$\hat\zeta$ & damping parameter, $c_rt_v/(2L)$; modal damping $\zeta_k=\hat\zeta\Omega_k$\\
$\tau$ & dimensionless time, $c_rt/L$\\
$\Omega=\omega L/c_r$ & dimensionless circular frequency\\
$\kappa_r$, $\kappa_c$ & dimensionless wavenumbers, $\Omega$ and $\gamma\Omega$\\
$W_k$, $Y_k$, $m_k$, $\Gamma_k$ & mode shape, take-up amplitude, modal mass and participation factor of mode $k$\\
$H_A$, $H_B$; $\tilde m$ & receptances of strands A and B at the take-up; modal mass of a strand mode seen from its free end\\
$\tper{1}$, $\tauper{1}$ & fundamental period of the loop; $\tauper{1}=c_r\tper{1}/L$\\
$\tper{s}$, $\tauper{s}$ & fundamental period of a fixed--free strand\\
$\tauB$ & transit time of the strand downstream of the take-up, $1-\xi+\gamma$\\
$\betafive$ & take-up mass ratio at which a natural period lengthens by 5\,\%\\
\bottomrule
\end{tabular}
\end{table}
